\documentclass[12pt,a4paper]{article}

% --- Pakiety podstawowe ---
\usepackage[utf8]{inputenc}
\usepackage[T1]{fontenc}
\usepackage[polish]{babel}
\usepackage{geometry}
\geometry{margin=2.5cm}
\usepackage{hyperref} % Linki interaktywne
\usepackage{tabularx} % Elastyczne tabele
\usepackage{enumitem} % Lepsze listy

% --- Konfiguracja linków ---
\hypersetup{
    colorlinks=true,
    linkcolor=blue,
    filecolor=magenta,      
    urlcolor=cyan,
    pdftitle={Dokumentacja Projektu - \texttt{SCUOLA}},
}

\begin{document}

% --- STRONA TYTUŁOWA ---
\begin{titlepage}
    \centering
    \vspace*{2cm}
    {\Huge\bfseries Projekt: System do zarządzania szkołą korepetycji \par}
    \vspace{1cm}
    {\Large Nazwa kodowa: \texttt{SCUOLA} \par}
    \vspace{2cm}
    
    {\large \textbf{Skład zespołu:}} \\
    \vspace{0.5cm}
    \begin{center}
        Bartosz Badura (Lider) \\
        Eryk Witkowski \\
    \end{center}
    
    \vfill
    
    {\large \textbf{Linki interaktywne:}} \\
    \vspace{0.5cm}
    \href{https://github.com/badubart/inzOprogramowania}{Repozytorium GitHub} \\
    \vspace{0.2cm}
    \href{https://trello.com/b/pA3IAemn/projekt-scuola}{Tablica trello} \\
    
    \vfill
    {\large 15 kwietnia 2026}
\end{titlepage}

% --- WIZJA SYSTEMU ---
\section{Wizja systemu}
\medskip

System \texttt{SCUOLA} ma na celu ułatwienie zarządzania szkołą korepetycji poprzez zautomatyzowanie procesów takich jak rejestracja uczniów, zarządzanie grafikami zajęć, oraz komunikację między nauczycielami a uczniami. System będzie miał warstwę desktopową, przeznaczoną dla administratorów i nauczycieli w systemie Windows, oraz warstwę webową, dostępną dla uczniów, umożliwiającą im przeglądanie dostępnych zajęć i rezerwację miejsc.

% --- ZAKRES USŁUG ---
\section{Zakres usług}

\subsection{Funkcjonalności planowane do realizacji}
\begin{itemize}
    \item Umożliwia uczniom przeglądanie dostępnych zajęć i rezerwację miejsc.
    \item Zapewnia administratorom i nauczycielom narzędzia do zarządzania grafikami zajęć i komunikacji z uczniami.
    \item Automatyzuje proces tworzenia harmonogramów zajęć na podstawie dostępności nauczycieli i uczniów.
    \item Zapewnia funkcjonalność powiadomień, informując uczniów o nadchodzących zajęciach i zmianach w harmonogramie.
    \item Umożliwia nauczycielom na bieżąco monitorować frekwencję uczniów i zarządzać listami obecności.
\end{itemize}

\subsection{Ograniczenia (Czego nie będziemy realizować)}
\begin{itemize}
    \item Integracja z zewnętrznymi systemami płatności.
    \item Funkcjonalność mobilna (aplikacja na smartfony).
    \item Zaawansowane funkcje analityczne i raportowanie.
    \item Obsługa wielu języków (system będzie dostępny tylko w języku polskim).
    \item Integracja z systemami operacyjnymi innymi niż Windows (np. macOS, Linux).
\end{itemize}

% --- ARCHITEKTURA LOGICZNA ---
\section{Architektura logiczna}
Niniejsza sekcja opisuje strukturę systemu oraz sposób komunikacji między jego elementami w nawiązaniu do środowiska operacyjnego.
\newpage
\begin{table}[h]
    \centering
    \renewcommand{\arraystretch}{1.5}
    \begin{tabularx}{\textwidth}{|l|X|X|}
    \hline
    \textbf{Element składowy} & \textbf{Zadania i odpowiedzialność} & \textbf{Komunikacja} \\ \hline
    Frontend (React) & Interfejs użytkownika, walidacja danych wejściowych. & Wywołania API REST do Backendu. \\ \hline
    Backend (Java/Spring Boot) & Logika biznesowa, autoryzacja, przetwarzanie danych. & SQL do bazy danych. \\ \hline
    Baza danych (MySQL) & Przechowywanie danych w sposób trwały. & Odpowiada na zapytania Backendu. \\ \hline
    \end{tabularx}
    \caption{Główne elementy składowe systemu i ich interakcje.}
\end{table}

% --- GŁÓWNE FUNKCJE (BACKLOG) ---
\section{Główne funkcje (Historyjki i Epiki)}

\subsection{Epik 1: Zarządzanie grafikiem zajęć}
\begin{description}
    \item[User Story 1:] Jako Uczeń, chcę wejść na stronę internetową szkoły i zobaczyć dostępne zajęcia, aby szybko zarezerwować lekcję.
    \item[User Story 2:] Jako uczeń chciałbym wybrać godziny w których będę uczęszczał na zajęcia.
    \item[User Story 3:] Jako uczeń chciałbym mieć możliwość podglądu dostępnych godzin i nauczycieli.
    \item[User Story 4:] Jako uczeń chciałbym móc wykonać rezerwację na dany okres czasu np. 2 miesiące.
    \item[User Story 5:] Jako uczeń chciałbym mieć możliwość zobaczenia dostępnych przedmiotów.
    \item[User Story 6:] Jako uczeń chciałbym mieć możliwość anulowania zajęć.
    \item[User Story 7:] Jako uczeń chciałbym widzieć cenę zajęć przed podjęciem decyzji.
    \item[User Story 8:] Jako uczeń chcę móc wyeksportować mój plan do kalendarza.
    \item[User Story 9:] Jako nauczyciel chciałbym widzieć swój plan.
    \item[User Story 10:] Jako nauczyciel chciałbym wybrać godziny, w których mogę pracować.
    \item[User Story 11:] Jako administrator chciałbym przypisać nauczyciela do sali.
    \item[User Story 12:] Jako administrator chciałbym aby system automatycznie zwalniał salę gdy ktoś anuluje zajęcia.
    \item[User Story 13:] Jako administrator chciałbym mieć podgląd do ogólnego planu zajęć.
    \item[User Story 14:] Jako administrator chcę mieć możliwość ręcznego wpisania nauczyciela i sali w nieprzewidzianej sytuacji.
    \item[User Story 15:] Jako administrator chciałbym aby system informował mnie o uczniach anulujących zbyt dużą ilość zajęć.
    \item[User Story 16:] Jako administrator chciałbym móc ręcznie przypisać ucznia gdyby on sam nie był w stanie.
\end{description}


\subsection{Epik 2: System powiadomień}
\begin{description}
    \item[User Story 1:] Jako Uczeń, chcę otrzymać potwierdzenie rezerwacji SMS, aby mieć pewność, że moja rezerwacja została przyjęta.
    \item[User Story 2:] Jako uczeń chciałbym otrzymać przypominającego SMS w dniu przed zajęciami.
    \item[User Story 3:] Jako nauczyciel chciałbym otrzymać SMS informującego o nowym uczniu.
    \item[User Story 4:] Jako nauczyciel chciałbym otrzymywać powiadomienia w przypadku anulowania zajęć.
    \item[User Story 5:] Jako administrator, chcę zobaczyć nową rezerwację z internetu.
    \item[User Story 6:] Jako administrator, chcę otrzymać powiadomienie w przypadku anulowania zajęć.
\end{description}
    
\subsection{Epik 3: Monitorowanie uczniów}
\begin{description}
	\item[User Story 1:] Jako uczeń, chcę mieć dostęp do udostępnionych przez nauczyciela notatek.
    \item[User Story 2:] Jako nauczyciel, chcę zalogować się do systemu na komputerze w sali, żeby odznaczyć, kto przyszedł na dzisiejszą lekcję.
    \item[User Story 3:] Jako nauczyciel chciałbym wyświetlić listę aktualnych uczniów.
    \item[User Story 4:] Jako nauczyciel chciałbym widzieć numer do ucznia/rodzica.
    \item[User Story 5:] Jako nauczyciel chciałbym mieć podgląd do obecności ucznia na zajęciach.
    \item[User Story 6:] Jako nauczyciel chciałbym zapisywać postępy danego ucznia.
    \item[User Story 7:] Jako nauczyciel chciałbym dodawać notatki do danego ucznia.
     \item[User Story 9:] Jako administrator chciałbym aby system informował mnie o uczniach, którzy opuścili zbyt dużą ilość zajęć.
    \item[User Story 8:] Jako manager chciałbym widzieć obecność wszystkich uczniów.
   
\end{description}




\end{document}